\documentclass[10pt,letterpaper]{article}
\usepackage[margin=23mm]{geometry}
\usepackage{amsmath,amssymb,booktabs,hyperref}
\hypersetup{colorlinks=true,urlcolor=blue,
  pdftitle={Separate-clock Gaussian reference for the full C. elegans lineage}}
\setlength{\parindent}{0pt}
\setlength{\parskip}{5pt}
\newcommand{\Normal}{\mathcal N}
\title{\textbf{Separate-clock Gaussian reference for the full
  \textit{C. elegans} lineage}}
\author{Embryogenesis Pareto analysis project}
\date{Adopted: 9 September 2026; reaffirmed: 16 September 2026}
\begin{document}
\maketitle

\section{Publication decision and scope}

Figure 8 uses a fixed-topology Gaussian reference with tracking-time scaling
for spatial changes and one unit per canonical lineage transition for protein
changes. This replaces the audited model that used tracking duration for all
23 features. The reason for adoption is compatibility between measurement
windows and the assumed clocks, not agreement with the observed total cost.
The earlier shared-time model remains available as a sensitivity comparison.

The final manuscript review retains the root-fixed, leaf-unconstrained
reference. Covariance is estimated from all internal and terminal edges;
forward simulation then generates all descendants without requiring the
observed leaf outcomes. Conditioning on those outcomes with the same fitted
covariance is mathematically valid, but asks a different question.

The configuration is unchanged: three spatial coordinates from tracking
embryo 1 through $T\leq255$, scaled by 0.1625; 20 selected z-scored protein
features from the aggregate atlas; four fixed measured roots; 500 internal and
504 terminal measured cells; and exactly 1,000 scored parent--child edges.
P0, AB and P1 are construction anchors and contribute no scored edges.
The selected features, preprocessing, natural costs, optimization strategies
and first-cousin-null plotting scales are unchanged by this decision.

\section{Why the protein clock was changed}

Tracking duration on an edge $e=(u\to v)$ is
$t_e=t^{\rm last}_v-t^{\rm last}_u$, after truncation at 255. Of the 504
terminal cells, 492 reach that cutoff. The protein states are separately
aggregated atlas measurements, not observations at these tracking endpoints.
There is therefore no basis for assuming that their squared differences
shrink in proportion to the remaining tracking interval.

In the former model the fitted covariance was
\[
 \widehat\Sigma_{\rm old}=\frac1{1000}\sum_e
   \frac{\Delta\mathbf x_e\Delta\mathbf x_e^\top}{t_e}.
\]
A large protein difference assigned a short interval contributes strongly
to this common per-time rate, which is then applied to every simulated edge.
The audit found that 49 of the 50 largest time-standardized protein changes
end at the cutoff, with median duration 5. The 118 edges with $t_e\leq10$
contribute 63.9\% of the fitted squared protein rate. The largest 5\% of
time-standardized increments contribute 56.0\% of that rate but only 10.8\%
of the observed sum of protein edge lengths.

For example, the edge \texttt{ABprpapppa} to \texttt{ABprpapppap} is assigned
tracking duration $255-250=5$, while its dominant reporter ZAG-1 observes
the daughter at 67 times on a different timeline. These raw reporter times
are provenance information, not replacement tracking durations. A cutoff
alone would not invalidate a Brownian model if the states were measured at
the corresponding endpoints; the concern is the mismatch of observation
windows. Equal protein variance per transition removes that unsupported
time scaling without introducing a fitted clock exponent or extra rate class.

\section{Model and parameter estimation}

For spatial state $\mathbf r_v\in\mathbb R^3$ and protein state
$\mathbf p_v\in\mathbb R^{20}$, define measured increments
$\Delta\mathbf r_e=\mathbf r_v-\mathbf r_u$ and
$\Delta\mathbf p_e=\mathbf p_v-\mathbf p_u$. The reference assumes
\begin{align}
 \Delta\mathbf r_e &\sim \Normal(\mathbf0,t_e\Sigma_r),\\
 \Delta\mathbf p_e &\sim \Normal(\mathbf0,\Sigma_p).
\end{align}
Increments are independent across edges, and the spatial and protein blocks
are independent. Within-block feature correlations are retained. Zero drift
is fixed; covariance fitting does not subtract a fitted mean.
The maximum-likelihood estimates are
\begin{align}
 \widehat\Sigma_r &= \frac1N\sum_{e=1}^N
     \frac{\Delta\mathbf r_e\Delta\mathbf r_e^\top}{t_e},\\
 \widehat\Sigma_p &= \frac1N\sum_{e=1}^N
     \Delta\mathbf p_e\Delta\mathbf p_e^\top,
 \qquad N=1000.
\end{align}
The stored block-diagonal covariance is
$\operatorname{diag}(\widehat\Sigma_r,\widehat\Sigma_p)$.
Its blocks have different units: squared spatial units per tracking-time
unit, and squared standardized protein units per canonical transition.
Every evaluated edge is checked to span exactly one canonical generation;
the implementation rejects a represented edge that skips generations.

This is Brownian scaling for position and a discrete Gaussian random walk
for protein state. It does not posit protein diffusion per minute. Setting
cross-block covariance to zero leaves each marginal cost distribution
unchanged when its within-block fit and clock are held fixed. Changing the
protein clock is the substantive modification; the two total costs need
not retain the former joint correlation.

\section{Simulation, scoring and reproducibility}

Fix only the four root states at their observed values; do not fix leaf states.
Traverse the measured
forest from roots toward tips. On each edge draw the two Gaussian increments
above and add them to the simulated parent state. The implementation validates
that reconstructed parent--child differences reproduce the drawn increments.
For each replicate score
\[
 C_r=\sum_{e\in E}\|\Delta\mathbf r_e\|_2,
 \qquad C_p=\sum_{e\in E}\|\Delta\mathbf p_e\|_2.
\]
These are the same edges and unsquared Euclidean norms used for the natural
lineage. No topological assignments are randomized. Although topology
determines shared ancestry of simulated states, these edge-additive costs
under independent increments depend on edge scales and covariance, not on
how the edges are connected. The reference is not a topology test.

The publication run uses NumPy's default random generator with seed 242,
10,000 replicates, batches of 100, and an eigendecomposition-based covariance
square root. The first 1,000 replicates form the deterministic display subset;
cloud means use all 10,000. Both covariance blocks must be positive definite.
Plotting coordinates subtract the natural cost and divide each objective by
its first-cousin-shuffle standard deviation. These plotting units are not
standard deviations of the Gaussian reference distribution.

\section{Results and remaining model inadequacy}

\begin{center}
\begin{tabular}{lrr}
\toprule
 & Travel & Protein state\\
\midrule
Observed & 4,465.55 & 3,102.41\\
Former shared-time mean & 4,831.32 & 5,240.78\\
Adopted separate-clock mean & 4,831.80 & 3,603.77\\
Adopted standard deviation & 72.97 & 25.51\\
Adopted central 95\% lower bound & 4,690.48 & 3,554.64\\
Adopted central 95\% upper bound & 4,977.13 & 3,653.72\\
\bottomrule
\end{tabular}
\end{center}

Protein excess falls from 68.9\% to 16.2\%. A second 10,000-replicate run
with seed 243 yields adopted means 4,831.11 and 3,603.97. Spatial marginal
behavior is unchanged apart from Monte Carlo variation.

The reference still poorly describes protein edge heterogeneity. Observed
edge-length coefficient of variation is 0.644 versus simulated mean 0.221;
the largest 5\% of raw squared increments contribute 30.6\% in the data
versus 11.3\% in simulations. Mean observed internal- and terminal-edge
lengths are 2.558 and 3.638, while the reference predicts about 3.604 for
both. Some source values were zero-imputed before z-scoring; three of the
top 50 historical rate outliers touch a missing selected-protein value.
Raw observation-window audits and missingness flags accompany the sensitivity
outputs, without changing preprocessing.

The cloud is a geometric parametric reference, not a realistic model of
embryogenesis. A smaller total-cost discrepancy alone does not validate the
Gaussian assumption. Neither its residual displacement nor descriptive tail
counts establish selection or biological efficiency. Plug-in simulations do
not include parameter uncertainty, measurement error, feature-selection
uncertainty or shared observation-error dependence between adjacent edges.
Any future branch-class or heavy-tailed extension should be motivated by
edgewise diagnostics and assessed using held-out subtrees before publication
adoption. This follows the distinction between absolute adequacy and relative
model fit described by \href{https://doi.org/10.1086/682022}{Pennell et al. (2015)}.

\section{Implementation and archived specifications}

Two reproducible side analyses were reviewed but not adopted for manuscript
figures. A lognormal protein branch multiplier gives mean protein cost about
2,988, but still misses seven of eight held-out subtree total intervals.
Joint root-and-leaf conditioning gives mean travel 4,066 and protein 2,936;
it reduces uncertainty and permits smoothed internal configurations without
developmental feasibility constraints. Lower or closer costs alone do not
establish adequacy or biological efficiency. The full specifications, results
and reproduction commands are in
\path{full_tree_pareto/BRANCH_VARIABILITY_SENSITIVITY.md} and
\path{full_tree_pareto/LEAF_CONDITIONED_REFERENCE.md}.
Neither changes the adopted reference or its publication caches.

The authoritative generator is \texttt{separate\_clock\_reference()} in
\path{full_tree_pareto/publication_analysis.py}. The Figure 8 renderer
requires the separate-clock model label and provenance. Covariance and
diagnostic files use the \path{ce_full_tree_separate_clock_} prefix;
the collective null and summary tables also identify the adopted model.
The companion \path{clock_sensitivity.py} compares both models, records
edge distributions and source windows, and generates diagnostic panels.

\path{parametric_brownian_bootstrap.tex} documents the superseded
shared-time model. \path{brownian_covariance_mle_derivation.tex}
remains a valid general covariance derivation: apply it to spatial increments
divided by $\sqrt{t_e}$ and separately to unscaled protein increments.
Neither the earlier full 23-dimensional shared-time model nor the incompatible
historical 10-PC ancestral-state cache should be reintroduced into Figure 8.
\end{document}
